% Created 2026-06-26 Fri 13:12
% Intended LaTeX compiler: pdflatex
\documentclass{article}
\usepackage[utf8]{inputenc}
\usepackage[T1]{fontenc}
\usepackage{graphicx}
\usepackage{longtable}
\usepackage{wrapfig}
\usepackage{rotating}
\usepackage[normalem]{ulem}
\usepackage{amsmath}
\usepackage{amssymb}
\usepackage{capt-of}
\usepackage{hyperref}

\usepackage{fancyhdr}
\usepackage[top=25mm, left=25mm, right=25mm]{geometry}
\usepackage{longtable}
\fancyhead[R]{}
\setlength\headheight{43.0pt}
\author{Almeida Anthonny, Haro Damian, Molina Jhanavi y Sagñay Micaela}
\date{2026-06-26}
\title{Taller de Comandos Linux}
\hypersetup{
 pdfauthor={Almeida Anthonny, Haro Damian, Molina Jhanavi y Sagñay Micaela},
 pdftitle={Taller de Comandos Linux},
 pdfkeywords={},
 pdfsubject={},
 pdfcreator={Emacs 27.1 (Org mode 9.7.5)}, 
 pdflang={Español}}
\begin{document}

\maketitle
\fancyhead[C]{\includegraphics[scale=0.05]{../images/logoEPN.jpg}\\
ESCUELA POLITÉCNICA NACIONAL\\FACULTAD DE INGENIERÍA DE SISTEMAS\\
ARQUITECTURA DE COMPUTADORES}
\thispagestyle{fancy}


\section{Objetivos}
\label{sec:org4fac5f8}

\begin{itemize}
\item Conocer el funcionamiento del terminal de \emph{Linux}
\item Aprender a manejar comandos del terminal de \emph{Linux}
\end{itemize}

\section{Instrucciones}
\label{sec:org501dca6}
Este taller/tutorial de comandos de Linux tiene por objetivo ejercitar
en la aplicación de comandos estudiados en clase. Para esto siga las
Instrucciones indicadas en cada uno de los numerales. Ejecute los
ejercicios llenando los comandos adecuados dentro de los bloques de
código de Emacs Org. Los resultados de la ejecución se obtienen
haciendo \texttt{C-c C-c}

Para todos los ejercicios se sugiere o bien usar directamente un
distro de \emph{Linux} o el \emph{WSL} configurado para la clase de
\emph{Arquitectura de Computadores}.

\begin{enumerate}
\item Utilice \texttt{cd} para localizarse en la posición de home de su usuario y
verifique ejecutando el comando \texttt{pwd}

\item Cree tres directorios denominados FashionData, Train y
Test. Verifique usando el comando \texttt{ls}

\item Investigue sobre el comando \texttt{wget} y utilícelo para descargar en
\textasciitilde{}/FashionData los datos de train y test del /Fashion-Mnist desde el
siguiente enlace: \href{https://github.com/zalandoresearch/fashion-mnist}{fashion-mnist}
\begin{enumerate}
\item Imágenes de entrenamiento:

\url{http://fashion-mnist.s3-website.eu-central-1.amazonaws.com/train-images-idx3-ubyte.gz}
\item Etiquetas de entrenamiento:

\url{http://fashion-mnist.s3-website.eu-central-1.amazonaws.com/train-labels-idx1-ubyte.gz}
\item Imágenes de prueba:

\url{http://fashion-mnist.s3-website.eu-central-1.amazonaws.com/t10k-images-idx3-ubyte.gz}
\item Etiquetas de prueba:

\url{http://fashion-mnist.s3-website.eu-central-1.amazonaws.com/t10k-labels-idx1-ubyte.gz}
\end{enumerate}

\item Utilice el comando \texttt{mv} de manera recursiva para pasar los datos de
entrenamiento a la carpeta \textasciitilde{}/Train

\item Utilice el comando \texttt{mv} de manera recursiva para pasar los datos de
prueba a la carpeta \textasciitilde{}/Test

\item Verifique que la carpeta FashionData está vacía

\item Los archivos anteriores están comprimidos. Investigue como
descomprimir usando el comando \texttt{tar}. Si es necesario realice la
instalación del comando. Descomprima los archivos en sus
respectivas carpetas. Apunte el comando utilizado para la descompresión.
\begin{verbatim}
   # Los archivos .gz de fashion-mnist no son tar.gz por lo tanto se usa gunzip en lugar de tar
   gunzip ~/Train/train-images-idx3-ubyte.gz
   gunzip ~/Train/train-labels-idx1-ubyte.gz
   gunzip ~/Test/t10k-images-idx3-ubyte.gz
   gunzip ~/Test/t10k-labels-idx1-ubyte.gz
\end{verbatim}

\begin{figure}[htbp]
\centering
\includegraphics[width=.9\linewidth]{/home/jhanavi/ArquiB2/ArquiB2-/images/imagen1.jpg}
\caption{Resultado ejecutado en terminal WSL}
\end{figure}
\begin{figure}[htbp]
\centering
\includegraphics[width=.9\linewidth]{/home/jhanavi/ArquiB2/ArquiB2-/images/imagen2.jpg}
\caption{Resultado ejecutado en terminal WSL}
\end{figure}

\item Una vez completa la descompresión, elimine los archivos
.gz. Verifique las carpetas finales usando \texttt{ls}
\begin{verbatim}
   # Como utilizamos gunzip, elimina automaticamente el .gz al descomprimir
   rm -f ~/Train/*.gz
   rm -f ~/Test/*.gz
\end{verbatim}

\begin{figure}[htbp]
\centering
\includegraphics[width=.9\linewidth]{/home/jhanavi/ArquiB2/ArquiB2-/images/imagen3.jpg}
\caption{Ejecutando en terminal WSL}
\end{figure}

\item Mueva recursivamente las carpetas \textasciitilde{}/Train y \textasciitilde{}/Test dentro de \textasciitilde{}/FashionData
\begin{verbatim}
   mv ~/Train ~/FashionData/ 
   mv ~/Test ~/FashionData/
\end{verbatim}

\begin{figure}[htbp]
\centering
\includegraphics[width=.9\linewidth]{/home/jhanavi/ArquiB2/ArquiB2-/images/imagen4.jpg}
\caption{Evidencia de que en FashionData se muestra Train y Test}
\end{figure}
\end{enumerate}



\begin{enumerate}
\item Se desea realizar un registro climatológico de la ciudad de
Quito. Para esto, escriba un script de Python/Java que permita
obtener datos climatológicos desde el API de
\href{https://openweathermap.org/current\#one}{openweathermap}. Considere la latitud y la longitud de Quito como
(-0.2299, -78.5249), respectivamente. Los resultados obtenidos de
la consulta al API se escriben en un archivo
\emph{clima-quito-hoy.csv}. Cada ejecución del script debe almacenar
nuevos datos en el archivo. Investigue sobre \textbf{crontab} y utilícelo
para obtener datos del API de \emph{openweathermap} cada 5 minutos
mediante la ejecución de un archivo ejecutable denominado
\emph{get-weather.sh}. Verifique los resultados. Todas las operaciones
se realizan en Linux o en el WSL. Las etapas del problema se
subdividen en:
\begin{enumerate}
\item Crear su API gratuito en \href{https://openweathermap.org/current\#one}{openweathermap}
\item Escribir un script en Python/Java que realice la consulta al API
y escriba los resultados en \emph{clima-quito-hoy.csv}

\item Desarrollar un ejecutable \emph{get-weather.sh} para ejecutar el
programa Python/Java.

\item Configurar Crontab para la adquisición de datos. Escriba el
comando configurado.
\end{enumerate}
\end{enumerate}
\end{document}
